%% ma: L12-B14-0080
%% lop: 12
%% chuong: 5
%% bai: 14
%% dang: 12.14.21
%% loai: TN4
%% muc: TH
%% dap_an: "B"
%% nguon: "To_n_12_PT_M_t_Ph_ng.pdf (Đề kiểm tra đánh giá kiến thức, Chuyên đề Phương trình mặt phẳng, Nguyễn Anh Khôi - Lê Minh Khôi), Câu 36"
%% kien_thuc: "Bài 14 (phương trình mặt phẳng theo đoạn chắn)"
%% trang_thai: cho-duyet
%% ly_do: "Dạng 12.14.21 mới đề xuất, chờ giáo viên duyệt."
%% ghi_chu: "Tài liệu gốc không kèm đáp án; đáp án do Claude giải và kiểm bằng sympy."
\begin{ex}
Trong không gian $Oxyz$, cho ba điểm $A(1;0;0)$, $B(0;-2;0)$, $C(0;0;-5)$. Vectơ nào dưới đây là một vectơ pháp tuyến của mặt phẳng $(ABC)$?
\choice{$\vec n_1=\left(1;\dfrac12;\dfrac15\right)$}{\True $\vec n_2=\left(1;-\dfrac12;-\dfrac15\right)$}{$\vec n_3=\left(1;-\dfrac12;\dfrac15\right)$}{$\vec n_4=\left(1;\dfrac12;-\dfrac15\right)$}
\loigiai{$(ABC)\colon \dfrac{x}{1}+\dfrac{y}{-2}+\dfrac{z}{-5}=1$ có vectơ pháp tuyến $\left(1;-\dfrac12;-\dfrac15\right)$. Chọn B.}
\end{ex}
